\documentclass[11pt]{article}
\usepackage{amsmath,amssymb}
\usepackage{graphicx}
\usepackage[hidelinks]{hyperref}
\usepackage{kotex}            % Korean draft — compile with XeLaTeX or LuaLaTeX
\usepackage[margin=1in]{geometry}

\title{관상동맥 파열위험지표 정의 및 민감도 분석}

\author{김세혁}

\date{}

\begin{document}

\maketitle

\begin{abstract}

관상동맥 죽상경화 플라크의 파열은 급성 관상동맥 증후군의 주요 기전으로서 임상적 중요성이 크며 이에 따른 예방과 즉각적 치료의 필요성은 진료지침에서도 강조된다 \cite{members2008}. 또한 바이오메카닉스 관점에서 플라크에 작용하는 기계적 힘은 병변의 진행과 파열 위험을 규정하는 핵심 인자이므로 힘학적 지표를 통한 위험 평가의 타당성이 제시되어 왔다 \cite{brown2016}. 그러나 기존 전산·영상 기반 연구들은 대체로 응력(stress)만을 파열위험의 척도로 삼아 환자마다 다른 섬유막의 강도(strength)를 반영하지 못했으므로, 파열 취약성은 응력을 환자 특이적 강도로 정규화한 비율 형태, 즉 $VI=\mathrm{stress}/\mathrm{strength}$ 로 정의되어야 더 정확한 위험평가가 가능하다.

이에 따라 플라크 파열을 결정하는 기전은 병변의 형태학적 특징(예: 목(shoulder)부의 기하학적 집중), 혈류역학적 부하(맥동성 압력·전단력) 및 조직의 재료강도(strength)가 상호작용하는 복합체계로 이해되어야 한다; 구조역학적 연구는 특히 circumferential stress 분포가 실제 파열 위치와 연관됨을 보여 플라크 내 응력 집중이 파열 기전에 직접적으로 관여함을 시사한다 \cite{cheng1993}, 그리고 바이오메카닉스 관점에서도 기계적 힘이 병변의 자연사에 중심적 역할을 한다는 점이 강조되어 왔다 \cite{brown2016}. 이 맥락에서 $VI=\mathrm{stress}/\mathrm{strength}$ 라는 응력-강도비 프레임워크는 재료가 분모로 작용하므로 환자별로 잘 알려지지 않거나 이질적인 섬유막 강도의 불확실성을 정량적으로 포함할 수 있는 합리적 물리적 근거를 제공한다 (구조역학적 정당성). 따라서 플라크 취약성 평가는 응력의 공간·시간적 분포뿐만 아니라 해당 부위의 실제 강도를 함께 고려하는 방식으로 재설계되어야 하며, 다음 절에서는 기존 연구들이 주로 응력만을 사용한 이유와 그 한계를 논의한다.

실제로 많은 전산적 분석들은 계산된 응력 분포를 파열 위험의 주요 지표로 사용해 왔으나, 측정 불확실성·계산·임상 획득의 제약으로 섬유막의 환자별 강도(strength)를 고정된 상수로 취급하거나 민감도 분석에서 배제하는 경우가 적지 않다; 이러한 접근은 응력 집중이 파열 위치와 연관된다는 관찰(\cite{cheng1993})에도 불구하고, 동일한 응력 수준에서 환자별 재료강도 차이에 따라 실제 취약성은 크게 달라질 수 있다는 중요한 불확실성을 남긴다. 따라서 응력만을 단독 지표로 사용하는 방법론적 관행은 파열 취약성의 환자 맞춤형 예측을 어렵게 하며, 앞서 정의한 취약성 지수(VI)가 재료 강도의 불확실성을 정량적으로 포함하도록 stress/strength 관점에서 재설계되어야 한다는 필요성을 낳는다. 이 같은 방법론적 격차는 어떤 인자군(형태·혈류역학·재료)이 실제 위험을 지배하는지를 밝히기 위한 광범위한 전역 민감도 분석의 중요성을 제기하며, 다음 절에서는 그 필요성을 구체적으로 논의한다.

앞서 지적한 응력·강도 불확실성을 고려하면, 형태학(morphology)·혈류역학(hemodynamics)·재료(materials) 세 인자군을 현실적 취약성지수(VI) 공간 위에서 동시에 넓게 변동시켜 전역적으로 분석하는 것이 필수적이다; 이는 단일 혈류 변수나 정적 응력 분포만으로는 포괄적 우선순위를 도출하기 어렵다는 점을 의미하며, 특히 서로 다른 혈역학적 지표들이 복합적으로 작용함을 지적한 선행 논의와도 맥을 같이 한다 \cite{meng2013}. 이러한 목적에는 입력 변수(또는 그룹)가 출력 분산에 기여하는 정도를 분해하여 주효과와 상호작용을 동시에 정량화하는 Sobol 1차($S_1$)·전차($S_T$) 지수가 적합하며, 이들 지수는 전역 민감도를 정량적으로 비교·정렬함으로써 민감도 기반 우선순위 도출의 통계적 근거를 제공한다. 다만, 맥동성 FSI 조건 하에서 세 인자군의 전체 설계공간을 고해상도로 샘플링하여 Sobol 지수를 안정적으로 추정하려면 막대한 수치계산 비용이 요구되므로, 기존 전산연구들은 주로 축약된 설정이나 제한적 변수군에서의 민감도 분석에 그치는 경우가 많았다; 따라서 본 연구에서는 이러한 계산적 제약을 극복하여 현실적 VI 위에서의 전역 민감도 분석을 수행하고자 한다.

이러한 계산적 제약을 극복하기 위해 본 연구에서는 선행 검증된 비용효율적 맥동 FSI 설정(경계조건·심근·유입 파형은 data/cost\_effective\_FSI\_BC/bc.md 및 data/input\_parameter/*.dat 파일을 참조)으로 두 표현형(LAP 3도메인, CP 4도메인)을 대상으로 형태·혈류역학·재료 3군을 1,000-sample 설계공간(data/input\_parameter/input.csv)에서 무작위 변동시켜 PSS·ΔPSS(구면평균) 및 파열위치 지수를 출력하는 대규모 데이터셋을 생성(M1)하였다; 생성된 출력으로부터 응력·강도 비 형태의 VI 후보 6개(T1)를 정의하고 임상적으로 확립된 7개 고위험 플라크 특징(T1)의 부호일치로 후보를 선별한 결과 최종적으로 VI1,$$\mathrm{VI}_1=\Delta\mathrm{PSS}/E_{FC}^{0.5},$$와 VI2,$$\mathrm{VI}_2=\Delta\mathrm{PSS}/E_{FC}^{1.0}$$를 채택하였다(C3.2, C3). 채택된 VI들에 대해 GPR surrogate를 학습(LAP 784 / CP 727 유효쌍)(T1)하고 surrogate 상에서 약 1.5만 회 평가에 해당하는 Sobol 1차($S_1$)·전차($S_T$) 전역 민감도 분석을 수행하여 민감도 분해를 얻었다(M2). 그 결과, 응력 관점에서는 ΔPSS(최대진폭, 피로 관점)가 PSS(최대응력)보다 더 강한 예측인자임이 관찰되어(ΔPSS 우수성; C2, C2.1) 맥동성(피로)을 반영한 해석의 필요성이 확인되었고, strength 정규화(지수 α의 선택)는 민감도와 지배 인자 해석에 실질적 영향을 미치는 것으로 나타났다(C3). 전역 민감도 분석에서 재료물성 그룹의 기여가 가장 컸으며(그룹 수준 Sobol 결과: LAP·VI2에서 Material $S_1\approx0.75$, 관련 파일 data/sobol/sobol\_grp\_LAP\_VI2.csv) 개별 파라미터 수준에서도 $E_{FC}$, $E_{vessel}$, SBP, PP가 상위 4인자로 확인되어($E_{FC}$ $S_1\approx0.55$, 관련 파일 data/sobol/sobol\_ind\_LAP\_VI2.csv) 재료물성>혈류역학>형태의 지배적 순서가 실증되었다(C4, C4.1, C4.2). 본 연구의 주요 기여는 비용효율적 맥동 FSI 기반 대규모 데이터·대리모형 프레임워크를 통해 ΔPSS 기반의 strength-정규화 VI가 임상적 고위험 특성과 잘 부합함을 제시하고, 동시에 전역 민감도 정량화를 통해 재료물성과 혈압 지표(SBP·PP)가 플라크 파열 취약성 변동을 주도함을 명확히 한 점이다(주요 정량적 근거는 상술한 Sobol 결과 파일들에 근거).

\end{abstract}

\section{Introduction}

임상적으로, 관상동맥 죽상경화 플라크의 파열은 급성 관상동맥 증후군(ACS)의 주된 원인이며 이로 인해 심혈관계 사망 부담이 상당함이 보고되어 왔다 \cite{members2008}; 이러한 임상적 중요성은 플라크 파열 위험을 정량적으로 평가할 필요성을 분명히 제기한다. 바이오메카닉스 관점에서도 응력과 변형 등 기계적 인자가 플라크의 자연사와 파열 메커니즘에서 핵심적 역할을 한다고 강조되어 왔으며, 이는 기계적 위험지표 개발의 이론적·임상적 동기를 제공한다 \cite{brown2016}.

이와 같은 임상·이론적 맥락에서 플라크의 파열 취약성은 단일 요인이 아닌 형태학적(플라크·관상동맥 기하), 혈류역학적(유동 및 맥동성 부하) 및 재료적(섬유막·기질의 강도와 피로 특성) 요인들이 상호작용하여 결정된다는 점이 강조된다; 특히 바이오메카닉스 관점에서 이러한 다인자 상호작용이 플라크의 자연사와 파열 위험을 규정한다는 논의가 제기되어 왔다 \cite{brown2016}, 또한 기계적 하중과 생물학적 약화가 복합적으로 작용하는 파열 메커니즘에 대한 고찰은 피로·응력 집중 및 조직 강도 저하의 상호연관성을 지적한다 \cite{arroyo1999}. 이러한 관점은 지오메트리와 유동이 국소 응력분포를 형성하고, 반복적 맥동 부하가 피로 손상을 누적시키며, 궁극적으로 재료적 약화가 파열에 기여한다는 복합적 경로를 시사하나, 어느 인자군이 우선적으로 지배적인지는 단일 연구로 일반화될 수 없으므로 모든 관련 인자를 함께 고려하는 평가틀이 필요하다.

이러한 맥락에서, 많은 전산적 연구들이 최대응력이나 국소응력과 같은 stress 계열 지표만으로 파열 취약성을 정의해 왔으나, 이들 연구의 다수는 환자마다 다른 섬유막·기질의 강도(strength)를 고정된 상수로 취급하거나 고려하지 않는 경우가 많다; 구조역학적 관점에서는 파열 가능성을 정량화할 때 응력과 강도를 직접 비교하는 응력‑강도 비(stress/strength)가 더 합리적인 물리적 프레임워크이다. 응력‑강도 비는 분모에 재료강도가 들어감으로써 환자별 강도의 이질성과 측정 불확실성을 정량적으로 반영할 수 있으며, 동일한 국소응력 분포에서도 강도 차이에 따라 취약성 판정이 달라질 수 있음을 의미한다. 또한 인간 관상동맥 병변에서의 원주(또는 원형) 응력 분포가 파열 위치와 관련됨을 보인 구조해석적·병리학적 선행연구가 있어 이 물리적 프레임워크의 근거를 제공한다 \cite{cheng1993}. 따라서 stress 기반 지표는 유용한 통찰을 주지만 환자 특이적 위험층화를 위해서는 strength 변동성을 포함한 stress/strength 관점의 도입이 필요하며, 이 점은 이후 어느 인자군이 실제 위험을 지배하는지 규명하기 위한 전역 민감도 분석의 필요성과 자연스럽게 연결된다.

이러한 배경에서, 형태·혈류역학·재료의 세 인자군을 strength를 반영한 현실적 취약성 지수(VI) 위에서 동시에 그리고 넓은 설계공간으로 변동시키며 어느 인자군이 파열 위험을 우선적으로 결정하는지를 규명하는 전역 민감도 분석이 필수적이다. 전역 민감도 분석은 단일 인자의 주효과뿐 아니라 인자들 간의 상호작용 기여까지 분해해야 하므로 통계적으로 견고한 분산 분해 기법이 필요하며, 이 목적에 부합하는 지표로서 Sobol 1차(S1) 및 전차(ST) 지수는 입력 파라미터(또는 그룹)가 출력 분산에 기여하는 정도를 분해하여 주효과와 상호작용을 정량화한다 — 이는 인자군별 우선순위를 데이터 기반으로 도출할 수 있는 통계적 근거를 제공한다 (W3). 다만, 세 인자군을 현실적 VI 범위에서 대규모로 표본화·평가하려면 전통적 수치해석에서 요구되는 계산 비용이 매우 크기 때문에 지금까지 이러한 포괄적 전역 민감도 분석을 수행한 연구는 제한적이었다; 따라서 본 연구에서는 다음 절에서 설명할 접근으로 이 격차를 메우고자 한다.

이를 위해 본 연구에서는 선행에서 검증된 비용효율적 맥동 FSI 파이프라인을 활용하여, 두 표현형(LAP 3도메인 및 CP 4도메인)에 대해 형태·혈류역학·재료의 3개 인자군을 포함하는 1,000-sample 설계공간을 광범위하게 표본화(입력 설계공간은 data/input\_parameter/input.csv 참조)하고 각 샘플에 대해 PSS, ΔPSS(구면평균) 및 파열위치 지수를 출력으로 생성하였다(경계조건 및 심근/유입 파형은 data/cost\_effective\_FSI\_BC/bc.md 및 관련 입력 파일들을 사용). 또한 본 연구는 stress 2종 × strength 3시나리오로 구성된 6개의 VI 후보를 정의한 뒤 임상적으로 정의된 7개 기준에 따라 후보를 선별하고, 최종적으로 채택된 VI에 대해 가우시안 과정 회귀(GPR) 기반 surrogate 모델을 학습(유효 쌍: LAP 784 / CP 727)하여 surrogate 위에서 Sobol 1차(S1) 및 전차(ST) 지수를 계산하는 전역 민감도 분석을 수행하였다(대략 1.5만 회 평가에 상응하는 surrogate 기반 반복 평가를 이용). 본 접근의 목적은 ΔPSS 기반의 stress/strength 취약성 지수들의 임상적 타당성을 평가하고, 현실적 입력 공간(강도 포함)에서 어떤 인자군이 파열 취약성을 지배하는지를 규명하는 데 있다.

\section{Methods}

본 연구의 분석 파이프라인은 데이터 생성 → 취약성 지수(VI) 정의·선별 → surrogate 기반 전역 민감도분석의 세 연속 단계로 구성된다. 먼저 데이터 생성 단계에서는 1,000개 샘플로 구성된 입력 설계공간(데이터 파일: data/input\_parameter/input.csv)을 기반으로 LAP·CP 등 표준 표현형에 대해 비용효율 맥동 FSI 해석을 수행하여 PSS, ΔPSS 및 파열위치 지표 등 필요한 출력변수를 산출한다. 다음으로 VI 정의·선별 단계에서는 취약성 지수를 $VI=\mathrm{stress}/\mathrm{strength}$로 규정하고, stress의 두 관점(예: 최대응력·최대진폭)과 strength의 세 시나리오를 조합하여 총 6개의 후보 VI를 구성한 뒤, 임상적으로 확립된 7개 고위험 기준에 대한 부호(일치) 검토를 통해 후보를 선별한다. 마지막으로 선택된 최종 VI에 대해 Gaussian process regression(GPR) 대리모델을 학습하고(LAP에서 784유효쌍, CP에서 727유효쌍 사용), 학습된 surrogate 위에서 Sobol 1차 및 전차(ST) 지수를 계산하여 전역 민감도분석을 수행한다(대략 1.5만 회에 해당하는 surrogate 평가에 상응). 다음 절(M-02)에서는 데이터 생성 단계의 수치적 설정(경계조건·입력파라미터 분포·시뮬레이션 절차)을 상세히 기술한다.

구체적으로 데이터 생성 단계에서는 선행에서 검증된 비용효율적 맥동 FSI 기법을 이용하여 두 표현형(LAP: 3도메인, CP: 4도메인)에 대해 data/input\_parameter/input.csv에 정의된 1,000-sample 설계공간(형태·혈류역학·재료 변수)을 무작위·라틴하이퍼큐브 등으로 샘플링한 입력 조합마다 맥동 FSI 해석을 수행하였다 \cite{terahara2020}. 경계조건 및 심근·유입 파형은 본 연구의 BC 파일을 직접 사용했으며(경계조건 상세: data/cost\_effective\_FSI\_BC/bc.md; 예: Rp=232.44 dyn·s/cm\textasciicircum 5, C=1.50E-03 cm\textasciicircum 5/dyn, Rd=1317.18 dyn·s/cm\textasciicircum 5, tcycle=1.0 s, CO=5.0 L/min, Emax=2.0 mmHg/cc 등), 심근 시간-탄성(elastance)과 관상동맥 유입파형은 각각 data/input\_parameter/elastance.dat 및 data/input\_parameter/inflow.dat 파일을 입력으로 사용하였다. 각 샘플에 대해 구조-유체 상호작용 해석으로부터 플라크 표면응력(PSS), 맥동성으로 인한 응력변동(ΔPSS; 구면평균 기준) 및 시뮬레이션 기반의 파열위치 지수를 계산하여 데이터셋으로 저장하였고, 이로써 후속 VI 정의·선별 및 surrogate 학습에 필요한 출력변수를 획득하였다. 다음 절(M-03)에서는 이렇게 생성된 출력으로부터 후보 취약성 지수의 수치적 계산 및 임상기준에 따른 선별 절차를 상세히 설명한다.

이어서, 데이터 생성 단계의 구체적 대상 모델과 가정은 다음과 같다. 본 연구에서는 두 가지 플라크 표현형을 명시적으로 고려하였는데, 저감쇠 플라크(LAP)는 3도메인(혈관벽·섬유막·지질/괴사핵)으로, 석회화 플라크(CP)는 4도메인(혈관벽·섬유막·지질/괴사핵·석회화)을 각각 분할하여 모델링하였다; 각 도메인의 탄성계수 등 재료 입력은 data/input\_parameter/input.csv의 컬럼(E\_vessel, E\_fc, E\_lipid, E\_cal 등)으로 지정되었다. 섬유막(fibrous cap, FC)은 모델 전반에서 두께가 공간적으로 균일하다고 가정하고 단일 스칼라 파라미터('fc\_av\_th' 컬럼)를 통해 샘플별로 변동시켰으며, 이 값은 모든 샘플에 대해 동일한 균일 두께로 할당되어 VI 계산시 강도(strength) 및 응력(stress) 정규화에 일관되게 사용되었다. 형태(예: DOS, lesion\_length, lumen\_axial\_skewness, lipid\_length\_ratio, d\_fc\_ca 등), 혈류역학(예: SBP, PP, DBP, PI, tau) 및 재료 파라미터는 data/input\_parameter/input.csv에 정의된 1,000-sample 설계공간에서 샘플링된 값을 그대로 적용하여 각 표현형별로 맥동 FSI 해석을 수행하였다. 경계조건 및 심근·유입 파형은 앞절에서 언급한 BC 파일(data/cost\_effective\_FSI\_BC/bc.md, data/input\_parameter/elastance.dat, data/input\_parameter/inflow.dat)을 사용하였다. 각 샘플·표현형 조합에 대해 얻은 출력(PSS, ΔPSS의 구면평균, 시뮬레이션 기반 파열위치 지수)은 다음 단계인 후보 VI의 수치 계산 및 임상기준 기반 선별에 입력으로 사용되며, 그 절차는 다음 절에서 상세히 기술한다.

이에 따라 본 절에서는 입력 설계공간을 명확히 규정한다: 본 연구는 형태(morphology), 혈류역학(hemodynamics), 재료(material)의 세 그룹으로 입력변수를 분류하고, 각 그룹의 파라미터를 조합한 1,000-sample 설계공간(data/input\_parameter/input.csv)을 사용하였다. 구체적으로 형태군 예시에는 DOS, lesion\_length, lumen\_axial\_skewness, lipid\_length\_ratio, d\_fc\_ca, fraction, ca\_axial\_skewness, ca\_shoulder\_skewness, fc\_av\_th 등이 포함되며; 혈류역학군에는 SBP, PP, DBP, PI, tau 등이 포함되고; 재료군에는 E\_vessel, E\_fc, E\_lipid, E\_cal 및 ca\_strength\_ratio 등 재료·강도 관련 변수가 포함된다(전체 변수와 샘플값은 data/input\_parameter/input.csv 참조). 설계는 전체 공간을 골고루 탐색하도록 라틴 하이퍼큐브 기반의 샘플링(또는 동등한 전역표집 전략)을 이용해 생성된 1,000개의 입력 조합을 사용하였고, 각 조합·표현형 쌍에 대해 경계조건 및 심근·유입 파형(data/cost\_effective\_FSI\_BC/bc.md, data/input\_parameter/elastance.dat, data/input\_parameter/inflow.dat)을 적용하여 PSS, ΔPSS(구면평균) 및 시뮬레이션 기반 파열위치 지수를 산출하였다. 이로써 후속 단계인 후보 VI의 수치계산 및 임상기준 기반 선별에 사용할 일련의 입력·출력 쌍이 확보되었으며, 다음 절(M-05)에서는 확보된 출력으로부터 6개 후보 VI를 수치적으로 계산하고 임상 기준에 따라 선별하는 절차를 상세히 서술한다.

앞서 기술한 입력·해석 절차에 따라, 각 샘플(및 표현형)에서 다음 세 종류의 출력변수를 표준화된 형식으로 기록하였다. 첫째, PSS(peak stress)는 해당 샘플의 맥동 FSI 시계열에서 플라크 표면에 관찰되는 최대 표면응력 값을 의미하며 각 샘플별로 스칼라 피크값으로 저장했다. 둘째, ΔPSS(응력 진폭; 구면평균)는 각 표면 지점에서의 시간적 응력 진폭(예: max−min)을 계산한 뒤 그 값을 표면 전체에 대해 구면평균(spherical mean)하여 얻은 스칼라로 정의·저장하였다(출력 변수명 및 입력은 data/input\_parameter/input.csv의 샘플 식별자와 대응시켜 기록). 셋째, 파열위치 지수는 해석 결과에서 최대응력 등 주요 응력특성의 공간적 분포를 요약한 스칼라 지표로서, 시뮬레이션 상의 위치 정보를 후속 분석에서 비교·정량화할 수 있도록 각 샘플마다 저장하였다. 이들 출력(PSS, ΔPSS(구면평균), 파열위치 지수)은 이후 절에서 설명할 후보 VI 수치계산 및 GPR surrogate 학습의 입력으로 사용되며, 다음 절(M-06)에서는 이 출력들로부터 6개 후보 VI를 실제로 계산하고 임상기준에 따라 선별하는 절차를 상세히 설명한다.

이어서, 후보 취약성 지수(VI)는 앞서 수집한 출력변수(PSS, ΔPSS)를 분자로 하고 섬유막 강도(strength)를 분모로 취하는 일반적 비율식으로 규정하여, stress(2종) × strength(3시나리오) 조합으로 총 6개의 후보를 구성하였다(C3.1). 구체적으로 분모인 strength에 대해서는 (A) 문헌상 관례와 같이 환자별 강도를 고정된 기준값으로 취하는 고정강도 시나리오(기호: S\_ref), (B) 섬유막 탄성계수 $E_{fc}$에 대해 지수 보정 $\alpha=0.5$를 적용하는 시나리오, (C) $E_{fc}$에 대해 $\alpha=1.0$을 적용하는 시나리오의 세 가지를 고려하였다. 일반식으로는 다음과 같이 쓸 수 있다(단, 고정강도 시나리오(A)에서는 분모를 상수 $S_{\mathrm{ref}}$로 표기한다): 
$$
VI\_{s,\textbackslash alpha}=\textbackslash frac{s}{E\_{fc}\textasciicircum {\textbackslash alpha}}\textbackslash qquad(\textbackslash alpha>0),
$$
$$
VI\_{s,{\textbackslash rm fixed}}=\textbackslash frac{s}{S\_{\textbackslash mathrm{ref}}}\textbackslash qquad({\textbackslash rm fixed\textbackslash  strength}),
$$
여기서 $s$는 stress 측정치(예: $s=\mathrm{PSS}$ 또는 $s=\Delta\mathrm{PSS}$)를 의미한다. 따라서 6개의 후보는 명명법으로 다음과 같다: VI\_{PSS,fixed}=PSS/S\_ref, VI\_{PSS,0.5}=PSS/E\_{fc}\textasciicircum {0.5}, VI\_{PSS,1.0}=PSS/E\_{fc}\textasciicircum {1.0}, VI\_{\textbackslash Delta PSS,fixed}=\textbackslash Delta PSS/S\_ref, VI\_{\textbackslash Delta PSS,0.5}=\textbackslash Delta PSS/E\_{fc}\textasciicircum {0.5}, VI\_{\textbackslash Delta PSS,1.0}=\textbackslash Delta PSS/E\_{fc}\textasciicircum {1.0}. 각 후보의 수식적 정의는 후보 정리 표에 요약되어 문서화하였다(표 참조: Table N — <정의된 6개 후보 VI 수식>). \textbf{[DATA NEEDED: S\_ref — 고정강도 시나리오에서 사용된 기준 강도 값(문헌·정규화 상수) 기재]}.

이어서 후보 VI의 선별은 임상적으로 확립된 '7개 고위험 플라크 특징'에 대한 부호(일치/불일치) 검사라는 단순·명시적 규칙을 적용하여 수행하였다: 각 후보 VI에 대해 해당 지표가 임상기준이 기대하는 방향(예: 고위험 표지와 양(또는 음)의 상관관계으로 해석되는 부호)을 만족하는지 여부를 일곱 기준 모두에 대해 판정하고, 모든 기준과 부호가 일치하는 후보만을 통과시켰다(선택 절차 및 결과 요약: Table N — <선택 절차 및 7개 임상기준 부호일치 결과>); 임상기준의 근거는 관련 병리·임상 문헌에 근거한다 \cite{wal1999}. 이 부호일치 규칙에 따라 6개 후보 가운데 두 개가 선택되었으며, 선택된 지수는 각각 
$VI_{1}=\Delta\mathrm{PSS}/E_{fc}^{0.5}$ 및 $VI_{2}=\Delta\mathrm{PSS}/E_{fc}^{1.0}$ 이다. 본 선별은 오직 임상기준과의 방향성(부호) 일치성만을 기준으로 한 절차임을 명확히 하며, 후보들의 정량적 성능 비교와 통계적 우수성 평가는 결과절에서 별도로 제시한다. 선정된 $VI_{1},VI_{2}$는 이후 절(M-08)에서 GPR 대리모델 학습 및 전역 민감도분석의 대상으로 사용된다.

위에서 선별된 $VI_{1},VI_{2}$에 대해서는 각 표현형별로 별도의 Gaussian process regression(GPR) 대리모델을 학습하였다(사용된 유효 학습쌍: LAP 784 유효쌍, CP 727 유효쌍). GPR 학습은 출력($VI_{1},VI_{2}$)을 목표로 입력 설계공간의 유효 샘플들에 대해 수행되었고[ cite:what\_I\_need — Gaussian process surrogate methodology for global sensitivity analysis ], 학습된 surrogate는 원래의 맥동 FSI 계산보다 계산비용이 현저히 낮아 약 1.5×10\textasciicircum 4(≈1.5만)회의 저비용 평가를 가능하게 하여 전역 민감도분석을 위한 충분한 표본수를 확보하였다. 준비된 surrogate 위에서 Sobol 1차 지수($S_{1}$) 및 전차 지수($S_{T}$)를 그룹 수준(형태/재료/혈류역학)과 개별 파라미터 수준에서 계산하였으며, 그룹별 결과는 data/sobol/sobol\_grp\_LAP\_VI1.csv, data/sobol/sobol\_grp\_LAP\_VI2.csv, data/sobol/sobol\_grp\_CP\_VI1.csv, data/sobol/sobol\_grp\_CP\_VI2.csv에, 개별 파라미터 수준 결과는 data/sobol/sobol\_ind\_LAP\_VI1.csv, data/sobol/sobol\_ind\_LAP\_VI2.csv, data/sobol/sobol\_ind\_CP\_VI1.csv, data/sobol/sobol\_ind\_CP\_VI2.csv에 각각 저장하였다. 여기서 Sobol의 $S_{1}$는 특정 입력변수가 출력 분산에 기여하는 주효과(단독 기여)를, $S_{T}$는 해당 입력의 주효과와 다른 입력과의 상호작용을 포함한 총기여를 분해·정량화하는 지표로서 전역적(분산 기반) 민감도 정량화를 제공하며(이로써 민감도 기반 우선순위 도출의 통계적 근거를 마련한다) 분석은 위 파일들에 기록된 지수들을 근거로 수행되었다. 마지막으로 명시하건대, 본 surrogate는 전역 민감도분석을 위한 계산적 에뮬레이터로 사용되었으며, 맥동 FSI의 검증·대체를 의도하지는 않는다.

\section{Results}

우선, 본 연구에서 구축한 데이터셋의 물리적 타당성을 점검한 결과, 본 논문에서 사용한 경계조건·입력값(데이터: data/cost\_effective\_FSI\_BC/bc.md, data/input\_parameter/inflow.dat, data/input\_parameter/elastance.dat 및 설계공간 표본인 data/input\_parameter/input.csv)에 따라 계산된 PSS, FFR 및 파열위치 지수의 공간적·정량적 경향이 선행 임상·구조해석 결과와 일관됨을 확인하였다; 예컨대 응력 분포와 파열 발생 위치의 관계는 기존 병변의 원주 응력 분포 관찰과 유사하게 나타나며 \cite{cheng1993} \textbf{[CITATION NEEDED: cite:wal1999]}, 혈류역학적 지표와 병변 특성의 상관성도 이전 보고들과 조화된다 \cite{eshtehardi2012} \textbf{[CITATION NEEDED: cite:brown2016]}. 또한 동일한 시뮬레이션 집합에서 계산한 비교 분석에서 ΔPSS(맥동에 의한 최대진폭)가 정적 최대 PSS에 비해 파열 취약성 지표와 더 강한 관련성을 보이는 것으로 관찰되어(본 섹션 후속문단에서 자세히 제시) 맥동성 해석의 필요성을 실증적으로 뒷받침한다. 다만 이러한 일치성은 본 데이터셋과 기존 지견 간의 경향 비교에 근거한 것으로, 환자 특이적 예측성이나 비용효율 FSI 기법 자체의 새로운 검증을 주장하는 것은 아니며 임상적 인과결론 역시 본 시뮬레이션 일치만으로 단정할 수 없음을 명확히 한다 (다음 단락으로 결과의 정량적 민감도·상관 분석을 제시함).

이러한 시뮬레이션 일치성에 이어, 본 연구에서는 먼저 stress 유형(정적 최대 PSS, 맥동에 의한 최대진폭 $\Delta PSS$) 2종과 strength 규격화의 3가지 시나리오를 조합하여 총 6개의 후보 취약성지수(VI)를 정의하였으며(각 후보의 수식적 형태는 일반형 $\mathrm{stress}/E_{fc}^{\alpha}$로 문서화됨; 세부식은 방법절에 기술), 이들 후보를 데이터셋(설계공간: data/input\_parameter/input.csv) 전역에서 평가한 뒤 임상적으로 확립된 7개 고위험 플라크 특징과의 부호 일치(sign‑matching) 기준으로 비교하였다. 그 결과 6개 중 두 개의 지수만이 7개 임상기준과 모두 부호가 일치하여 최종적으로 VI1 = $\Delta PSS / E_{fc}^{0.5}$ 및 VI2 = $\Delta PSS / E_{fc}^{1.0}$을 채택하였다(선택 근거: 7개 임상 특징과의 부호 일치).

이어서, 선택된 VI들(VI1, VI2)에 대해 surrogate 기반 전역 민감도 분석(Sobol)을 수행한 결과 그룹 수준에서 재료(Material) 그룹의 기여가 형태(Morphology)·혈류역학(Hemodynamics)보다 현저히 컸다: LAP·VI2에서 Material 그룹의 1차 민감도 $S_1=0.7506966098655802$(data/sobol/sobol\_grp\_LAP\_VI2.csv), LAP·VI1에서 $S_1=0.501286013821807$(data/sobol/sobol\_grp\_LAP\_VI1.csv), CP·VI1에서 $S_1=0.5087122894884933$(data/sobol/sobol\_grp\_CP\_VI1.csv), CP·VI2에서 $S_1=0.7967778987247732$(data/sobol/sobol\_grp\_CP\_VI2.csv)로 관찰되어 전반적으로 Material > Hemodynamics > Morphology의 우위가 확인되었다 (그룹별 결과 요약: data/sobol/*\_grp\_*.csv; (Fig F1 — <group‑level Sobol indices>) 참조). 개별 파라미터 수준에서도 재료물성 및 혈압 지표가 상위 결정인자로 나타났는데, 특히 LAP·VI2에서 $E_{fc}$의 $S_1=0.5503271567987186$로 1위였고 상위 4인은 $E_{vessel}$, $E_{fc}$, SBP, PP로 집계되었다(data/sobol/sobol\_ind\_LAP\_VI2.csv; (Fig F2 — <parameter‑level Sobol indices>, T2)). 다른 경우들도 일관되어 예컨대 LAP·VI1에서는 $E_{vessel}$($S_1=0.3068211191558551$)·SBP($S_1=0.16617475822426872$)가 주요 기여를 보였고(data/sobol/sobol\_ind\_LAP\_VI1.csv), CP·VI2에서는 $E_{fc}$가 더욱 뚜렷히 우세($S_1=0.6272396200605556$)하였다(data/sobol/sobol\_ind\_CP\_VI2.csv); CP·VI1의 개별 결과도 data/sobol/sobol\_ind\_CP\_VI1.csv에 수록되어 있다. 다만 이들 Sobol 지수는 본 논문에서 정의한 VI와 본 연구의 설계공간(입력 샘플 범위: data/input\_parameter/input.csv) 내 출력 분산에 대한 기여도를 정량화한 것이며, 환자 수준의 인과성 증거로 해석해서는 안 되고(또한 샘플링 범위나 VI 정의가 달라지면 순위가 변할 수 있음), 임상적 권고나 일반화는 별도 검증을 필요로 한다.

\section{Discussion}

앞서 기술한 결과에 기반하여, 본 연구의 시뮬레이션·대리모델 기반 분석에서는 응력의 최대진폭인 ΔPSS(피로·진폭 관점)가 단일 시점의 최대응력(PSS)보다 파열 취약성 지표로서 더 일관된 설명력을 보였고, 실제로 6개 후보 VI 중 ΔPSS 기반 지표들이 임상적 고위험 특성들과의 부호 일치 및 선별 단계에서 우수하게 선택되었다(시뮬레이션 타당성 및 VI 선별 절차 참조; E\_sim\_validation). 이러한 관찰은 조직 수준에서 반복적 맥동에 의한 피로·열화가 파열에 기여한다는 기계적·병태생리학적 논의와 일치하며, 피로(응력 진폭)의 중요성은 기계적 관점의 총체적 검토와 기초적 병리 보고서들에서 지지된다 \cite{arroyo1999} \cite{cheng1993}. 따라서 파열 취약성 평가에서 응력 진폭을 정확히 포착하려면 맥동성(시간 변동)과 유체–구조 상호작용을 반영하는 해석이 필수적이며, 정적·단일시점 해석만으로는 ΔPSS가 제공하는 피로 관련 정보를 얻기 어렵다. 다만 본 결과는 모델 기반 근거이며, 임상 환자 집단에서의 인과성 또는 예측성 입증을 주장하는 것은 아니고, 또한 특정 상황에서는 peak PSS가 여전히 유의한 기계적 위험 지표일 수 있음을 명확히 한다. 다음으로는 strength 정규화의 필요성과 정규화 지수의 지수 α가 민감도 결과에 미치는 영향을 논의한다.

따라서 파열 취약성 평가는 단순한 응력 지표가 아니라 응력과 환자별 강도의 비, 즉 stress/strength 형태로 정규화되어야 한다는 논리적·구조역학적 근거가 있다; 재료가 분모로 작용함으로써 환자 간 강도 불확실성을 정량적으로 반영할 수 있기 때문이다(기초적 병태·기계적 논의 참조) \cite{arroyo1999} \cite{cheng1993}. 본 연구에서는 이러한 관점을 반영하여 stress(최대응력·최대진폭) 2종과 strength의 3가지 시나리오를 조합하여 총 2×3=6개의 후보 VI를 구성하였고(E\_vi\_candidates), 이들에 대해 임상적으로 확립된 7개 고위험 플라크 특징과의 부호 일치성을 기준으로 후보를 선별하여 VI1 = $\Delta\mathrm{PSS}/E_{FC}^{0.5}$ 및 VI2 = $\Delta\mathrm{PSS}/E_{FC}^{1.0}$를 최종 채택하였다(E\_vi\_selection). 이때 정규화 지수 α의 선택은 단순한 스케일링을 넘어서 VI의 민감도 분해와 지배 인자 해석을 변화시킨다: α가 작을수록 재료물성의 영향이 상대적으로 완화되고, α가 클수록 재료가 더 큰 분산 기여를 하여 재료군의 우세성을 강화할 수 있다. 다만 우리는 여기서 모든 환자에 대해 보편적으로 최적인 단일 α를 제시하지 않으며, strength 정규화가 환자 수준의 모든 불확실성을 완전히 해소한다고 주장하지 않는다; 대신 정규화된 VI와 α의 변동성이 민감도 분석 결과(다음 절)를 통해 어떻게 재료·혈류역학·형태의 영향력을 재분배하는지 보여주고자 한다.

다음으로, 전역 민감도 분석 결과를 토대로 파열 위험도 결정요인을 해석하면 재료물성(Material)이 혈류역학(Hemo)과 형태(Morpho)를 능가하여 우세한 기여를 보였고, 이로부터 Material > Hemo > Morpho의 순서가 도출된다. 그룹 수준의 Sobol 1차 지수(S1)는 LAP·CP 양 표현형에서 재료군이 가장 큰 기여를 보였으며(관련 원시값은 data/sobol/sobol\_grp\_LAP\_VI2.csv, data/sobol/sobol\_grp\_LAP\_VI1.csv, data/sobol/sobol\_grp\_CP\_VI1.csv, data/sobol/sobol\_grp\_CP\_VI2.csv 참조), 특히 LAP·VI2에서는 Material 그룹의 S1이 약 0.75로 관찰되었다(data/sobol/sobol\_grp\_LAP\_VI2.csv) — 이는 그룹 차원에서 재료 변동이 출력 분산의 주된 원천임을 의미한다(C4.1). 개별 파라미터 수준의 분해에서도 E\_FC(E\_fc), E\_vessel, 수축기혈압(SBP), 맥압(PP)이 상위 4인자로 나타났고(예: E\_FC의 S1 ≈ 0.55, data/sobol/sobol\_ind\_LAP\_VI2.csv), 이들 네 변수의 우선순위는 VI1·VI2 모두에서 일관되었다(C4.2; 관련 결과는 (Fig 1 — <group Sobol 요약>) 및 (Fig 2 — <개별 파라미터 Sobol>)와 (Table 2 — <민감도 수치 요약>)에 정리되어 있다). 이러한 해석의 통계적 근거는 Sobol 1차(S1) 및 전차(ST) 지수가 입력 변수(그룹)가 출력 분산에 기여하는 주효과와 상호작용 기여를 정량적으로 분해한다는 점에 있으며(W3), 본 연구에서는 S1·ST의 비교를 통해 재료물성의 직접효과와 상호작용을 모두 평가하였다. 임상적 함의로는, 본 모델 기반 결과가 재료물성—특히 관벽·섬유막 강성—및 혈압 특성의 정량적 평가가 위험층화에서 높은 정보가치를 가질 수 있음을 시사하지만, 이는 모델 기반 우선순위의 제안일 뿐 즉각적인 임상 지침이나 절대적 진단임계값을 제시하는 근거로 삼을 수는 없다; 또한 형태 관련 측정은 특정 환자나 병변에서 여전히 유의미한 정보를 제공할 수 있음을 유지해야 한다.

한편, 본 연구에는 해석적·모델링 관점에서 몇 가지 명확한 한계가 있다. 우선 데이터 생성 단계(M1)에서 사용한 비용효율적 맥동 FSI 설정(경계조건·심근·유입 파형 등)은 data/cost\_effective\_FSI\_BC/bc.md 및 data/input\_parameter/input.csv에 기술된 절차를 따랐고, 입력 설계공간(형태·혈류역학·재료)은 data/input\_parameter/input.csv에 기초한 1,000-sample 샘플링으로 구성되었지만(두 표현형: LAP 3도메인, CP 4도메인), 해석 모델은 이상화된 혈관 형상과 균일한 fibrous cap 두께(단일 스칼라 변수로 취급) 등의 가정을 포함하므로 실제 환자별 해부학적 다양성과 국소적 두께·구조 불균일성을 완전하게 반영하지 못한다. 둘째, VI 후보 정의·선별과 surrogate 기반 민감도 분석(M2)은 6개 후보 VI를 임상 7기준으로 선별한 뒤 최종 VI에 대해 GPR surrogate를 학습(LAP 784 / CP 727 유효쌍)하고 해당 surrogate 위에서 Sobol 1차·전차 지수를 산출하는 절차를 따랐으나, surrogate는 근사 모델이므로 근사오차와 학습 데이터의 분포 의존성으로 인해 실제 고비용 FSI의 출력 분포를 완벽히 대체하지 못할 수 있으며, surrogate 기반 Sobol 추정에는 이러한 근사불확실성이 반영되어야 한다(본 연구에서는 surrogate를 약 1.5만 회 평가에 활용하여 전역 민감도를 실용적으로 추정함). 마지막으로, 비용효율적 FSI 방법 자체의 상세한 검증·성능비교는 본 연구의 범위 밖이며 해당 기법의 검증은 선행 연구에 위임한다(관련 검증문헌 참조 필요) \textbf{[DATA NEEDED: cost-effective FSI validation citation]}. 따라서 본 연구 결과는 주로 모델 기반 우선순위와 메커니즘적 통찰을 제공하는 것이며, 모든 환자 해부학으로의 일반화나 비용효율적 FSI 기법의 신규 검증·우월성 주장으로 해석되어서는 안 된다; 추가적으로 환자특이형상, 국소 섬유막 강도 측정 및 고해상도 검증이 향후 연구에서 요구된다.

\section{Conclusion}

요약하면, 본 연구는 시뮬레이션 기반 평가에서 ΔPSS(응력 최대진폭, 피로 관점)를 기반으로 한 취약성지수(VI)가 PSS(최대응력) 기반 지표보다 임상적으로 확립된 고위험 플라크 특성과의 정합성 및 모델 내 설명력 측면에서 우수함을 확인하여(예: VI 선별 결과 비교 참조; Fig 1 — <VI 후보 성능 비교>, Table 2 — <VI와 임상기준 부호일치>), 최종적으로 ΔPSS 기반 지수(VI1/VI2)를 채택하였다. 전역 민감도 분석에서는 재료물성(Material) 그룹이 혈류역학(Hemo)·형태(Morpho)를 능가하는 지배적 기여를 보였으며(LAP·CP 양상에서 그룹 수준 Sobol 1차값에서 Material이 최대 기여를 차지함; 관련 원시 결과 파일: data/sobol/sobol\_grp\_LAP\_VI2.csv, data/sobol/sobol\_grp\_LAP\_VI1.csv, data/sobol/sobol\_grp\_CP\_VI1.csv), 특히 LAP·VI2의 그룹 수준에서 Material의 기여가 $S_1\approx0.75$로 관찰되었다(data/sobol/sobol\_grp\_LAP\_VI2.csv). 개별 파라미터 수준의 Sobol 분석에서는 섬유막 강성 및 혈관 탄성 등이 우세하였고(data/sobol/sobol\_ind\_LAP\_VI2.csv), 상위 4개 결정인자로는 $E_{vessel}$, $E_{FC}$, 수축기혈압(SBP), 맥압(PP)이 도출되었으며(예: $E_{FC}$의 $S_1\approx0.55$), 이는 재료물성과 혈압 관련 변수들이 본 모델 내에서 파열 취약성 변동을 주도함을 시사한다(관련 그래프/수치: Fig 2 — <Sobol 민감도 요약>, Table 2 — <개별 파라미터 Sobol 지수>). 다음 절에서는 이러한 정량적 결과의 임상적·해석적 함의와 한계를 논의한다.

이러한 정량적 결과를 임상적·실무적 우선순위로 환원하면, 본 모델·VI 조건에서 관찰된 바와 같이 재료물성 그룹이 위험 변동을 지배하는 경우(예: data/sobol/sobol\_grp\_LAP\_VI2.csv에서 Material의 $S_1\approx0.75$, 관련 요약: Fig 2 — <Sobol 민감도 요약>, 또한 data/sobol/sobol\_grp\_LAP\_VI1.csv 및 data/sobol/sobol\_grp\_CP\_VI2.csv가 유사한 경향을 보임)를 근거로 재료 특성의 환자별 측정·정량화를 우선적으로 강화하는 것이 합리적이다. 그 근거는 응력-강도비(VI)가 강도를 분모로 포함하므로(stress/strength 프레임워크), 환자 간 강도 불확실성이 VI의 불확실성으로 직접 전이되며 따라서 섬유막·혈관 탄성계수 등의 직접 측정 또는 신뢰도 높은 추정이 VI 신뢰성 향상에 핵심적이라는 구조역학적 논리(W1)에 있다. 실무적 예로는 플라크·혈관의 탄성 계수 표준화된 측정(예: 임상적 탄성화상법·역학적 역추정 프로토콜)과 동시에 수축기혈압(SBP)·맥압(PP) 같은 혈압 지표의 정밀 측정을 우선 적용하는 것이 타당하다. 다만 본 권고는 본 연구의 모델·VI·파라미터 공간에서의 민감도 결과를 바탕으로 한 실무적 우선순위 제안일 뿐이며, 재료 특성 측정이 임상 결과를 확정적으로 개선한다고 주장하지 않으며, 특정 상황에서는 형태 기반 영상이 여전히 중요한 정보를 제공할 수 있음을 배제하지 않는다.

따라서 본 연구의 결과는 몇 가지 명확한 방법론적 시사점을 제시한다. 첫째, 플라크 파열 취약성 평가는 단순 최대응력(PSS)보다 응력의 최대진폭(ΔPSS, 피로 관점)을 기본 스트레스 지표로 삼는 것이 임상적으로 더 적합하다는 결론을 뒷받침하며(본 연구의 ΔPSS 우수성 근거: Fig 1 — <VI 후보 성능 비교>, 관련 시뮬레이션 타당성 T1), 이는 파형 기반의 맥동성 해석이 필수적임을 의미한다(본 분석은 맥동 유입·심근 elastance 경계조건을 사용했음: data/input\_parameter/inflow.dat, data/input\_parameter/elastance.dat; 경계조건·설정 상세: data/cost\_effective\_FSI\_BC/bc.md). 둘째, 취약성지수는 근본적으로 stress/strength 형태로 정의되어야 하며 이를 수학적으로 쓰면 $$\mathrm{VI}=\frac{\mathrm{stress}}{\mathrm{strength}}$$이고 본 연구에서처럼 피로 관점의 stress를 사용하면 예를 들어 $$\mathrm{VI}=\frac{\Delta\mathrm{PSS}}{E_{FC}^{\alpha}}$$와 같은 정규화가 필요하다; 여기서 설계공간에 정규화 지수 $\alpha$를 포함시켜 변동시켰음(data/input\_parameter/input.csv의 'alpha' 열)으로써 $\alpha$의 선택이 민감도와 지배 인자 해석에 실질적 영향을 미친다는 점을 확인하였다. 셋째, 따라서 향후 연구·임상 적용에서는 맥동성(피로) 처리 능력이 있는 FSI 기반 해석을 표준으로 채택하고, strength 정규화(지수 $\alpha$ 포함)의 민감도 검증과 환자별 strength 추정 절차를 병행해야 하며(기계적 파열 메커니즘 관련 선행근거: \cite{cheng1993}, \cite{arroyo1999}), 단일 $\alpha$ 값이 보편적으로 옳다고 주장해서는 안 되므로($\alpha$ 선택은 추가 검증·임상 연계 연구가 필요함), 다양한 $\alpha$ 설정 하에서의 강건성 검토와 환자 특이적 강도 측정의 통합을 권고한다.

앞서 논의한 대로, 본 연구가 제시한 결론들은 중요하지만 몇 가지 한계와 후속 과제를 낳는다. 첫째, 취약성지수의 정의에 포함한 정규화 지수 $\alpha$에 대한 민감도가 결과 해석을 실질적으로 바꾸므로(본 연구는 input.csv의 'alpha' 열을 통해 $\alpha$를 설계변수로 포함시켰음) 향후 연구에서는 $\alpha$의 범위를 체계적으로 탐색하고 각 $\alpha$ 값에 따른 Sobol 민감도 변화를 정량화해야 한다(참고 가능한 원시 민감도 파일: data/sobol/sobol\_ind\_LAP\_VI2.csv, data/sobol/sobol\_ind\_LAP\_VI1.csv, data/sobol/sobol\_ind\_CP\_VI1.csv). 둘째, 민감도 분석에서 $E_{FC}$·$E_{vessel}$ 등 재료 파라미터가 우세하게 나타났다는 사실(data/sobol/sobol\_ind\_LAP\_VI2.csv; 관련 요약: Table 2 — <개별 파라미터 Sobol 지수>)은 환자별 재료물성의 직접 측정 또는 신뢰도 높은 추정이 VI 신뢰성 향상에 핵심임을 시사하지만, 본 연구에서는 환자 특이적 물성 데이터가 포함되지 않았으므로 향후에는 조직 역학적 실험·영상 기반 탄성 추정(예: 임상 탄성화상법 또는 역문제 기반 추정)과 본 VI를 결합한 검증 작업이 필요하다. 셋째, 제안한 VI(예: $\mathrm{VI}=\Delta\mathrm{PSS}/E_{FC}^{\alpha}$)의 임상적 유효성은 아직 전향적 환자코호트 또는 병리·임상 결과와의 직접 비교를 통해 검증되어야 하며, 본 연구는 그러한 임상 검증을 수행하지 않았음을 분명히 한다(관련 VI 선별 근거: Table 2 — <VI와 임상기준 부호일치>). 따라서 권장되는 구체적 다음 단계는 (1) 다양한 $\alpha$ 설정에 대한 민감도-강건성 스터디 수행, (2) 환자별 섬유막·혈관 탄성 계수 획득을 위한 체계적 측정·역추정 프로토콜 개발 및 소규모 검증 데이터 확보, (3) surrogate 기반 대규모 전산 실험과 병용한 전향적 임상 연계 검증을 통해 VI의 예측성·임상적 유용성을 평가하는 것이다. 이러한 단계들이 선행될 때까지 본 연구에서 제안한 지표들은 개념적·수치적 근거를 제공하는 단계에 머물며, 즉시 임상 배포를 주장할 수 없음을 재차 강조한다.

\begin{thebibliography}{99}
\bibitem{members2008} Authors/Task Force Members, Frans Van de Werf, Jeroen J. Bax, Amadeo Betriu, C. Blomström‐Lundqvist, Filippo Crea, Volkmar Falk, Gerasimos Filippatos, Keith A.A. Fox, Kurt Huber, Adnan Kastrati, Annika Rosengren, Philippe Gabríel Steg, Marco Tubaro, Freek W.A. Verheugt, Franz Weidinger, Michael Weis, Alec Vahanian, A. John Camm, Raffaele De Caterina, Verónica Dean, Kenneth Dickstein, Gerasimos Filippatos, Christian Funck‐Brentano, Irene Hellemans, Steen Dalby Kristensen, Keith McGregor, Udo Sechtem, Sigmund Silber, Michał Tendera, Petr Widimský, José Luis Zamorano, Document Reviewers, Sigmund Silber, Frank V. Aguirre, Nawwar Al‐Attar, Eduardo Alegría, Felicita Andreotti, Werner Benzer, Ole A. Breithardt, N Danchin, Carlo Di Mario, Dariusz Dudek, Dietrich C. Gulba, Sigrun Halvorsen, Philipp A. Kaufmann, Ran Kornowski, Gregory Y.H. Lip, Frans H. Rutten (2008). Management of acute myocardial infarction in patients presenting with persistent ST-segment elevation. \href{https://doi.org/10.1093/eurheartj/ehn416}{DOI: 10.1093/eurheartj/ehn416}
\bibitem{vandewerf2003} F VandeWerf (2003). Management of acute myocardial infarction in patients presenting with ST-segment elevation. \href{https://doi.org/10.1016/s0195-668x(02)00618-8}{DOI: 10.1016/s0195-668x(02)00618-8}
\bibitem{martntimn2014} Iciar Martín-Timón (2014). Type 2 diabetes and cardiovascular disease: Have all risk factors the same strength?. \href{https://doi.org/10.4239/wjd.v5.i4.444}{DOI: 10.4239/wjd.v5.i4.444}
\bibitem{meng2013} Hui Meng, Vincent M. Tutino, J. Xiang, Adnan H. Siddiqui (2013). High WSS or Low WSS? Complex Interactions of Hemodynamics with Intracranial Aneurysm Initiation, Growth, and Rupture: Toward a Unifying Hypothesis. \href{https://doi.org/10.3174/ajnr.a3558}{DOI: 10.3174/ajnr.a3558}
\bibitem{cheng1993} Guijuan Cheng, Howard M. Loree, Roger D. Kamm, Michael C. Fishbein, R T Lee (1993). Distribution of circumferential stress in ruptured and stable atherosclerotic lesions. A structural analysis with histopathological correlation.. \href{https://doi.org/10.1161/01.cir.87.4.1179}{DOI: 10.1161/01.cir.87.4.1179}
\bibitem{wal1999} Allard C. van der Wal (1999). Atherosclerotic plaque rupture – pathologic basis of plaque stability and instability. \href{https://doi.org/10.1016/s0008-6363(98)00276-4}{DOI: 10.1016/s0008-6363(98)00276-4}
\bibitem{brown2016} Adam J. Brown, Zhongzhao Teng, Paul C. Evans, Jonathan H. Gillard, Habib Samady, Martin R. Bennett (2016). Role of biomechanical forces in the natural history of coronary atherosclerosis. \href{https://doi.org/10.1038/nrcardio.2015.203}{DOI: 10.1038/nrcardio.2015.203}
\bibitem{arroyo1999} Luis G. Arroyo (1999). Mechanisms of plaque rupture mechanical and biologic interactions. \href{https://doi.org/10.1016/s0008-6363(98)00308-3}{DOI: 10.1016/s0008-6363(98)00308-3}
\bibitem{eshtehardi2012} Parham Eshtehardi, Michael McDaniel, Jin Suo, Saurabh S. Dhawan, Lucas H. Timmins, José Binongo, Lucas Golub, Michel Corban, Aloke V. Finn, John N. Oshinsk, Arshed A. Quyyumi, Don P. Giddens, Habib Samady (2012). Association of Coronary Wall Shear Stress With Atherosclerotic Plaque Burden, Composition, and Distribution in Patients With Coronary Artery Disease. \href{https://doi.org/10.1161/jaha.112.002543}{DOI: 10.1161/jaha.112.002543}
\bibitem{cullen2003} Paul Cullen, Roberta Baetta, Stefano Bellosta, Franco Bernini, Giulia Chinetti, Andrea Cignarella, Arnold von Eckardstein, Andrew Exley, Martin Goddard, Marten H. Hofker, Eva Hurt‐Camejo, Edwin Kanters, Petri T. Kovanen, Stefan Lorkowski, William L. McPheat, Markku O. Pentikäinen, Jürgen Rauterberg, Andrew J. Ritchie, Bart Staels, Benedikt Weitkamp, Menno P.J. de Winther (2003). Rupture of the Atherosclerotic Plaque. \href{https://doi.org/10.1161/01.atv.0000060200.73623.f8}{DOI: 10.1161/01.atv.0000060200.73623.f8}
\bibitem{terahara2020} Takuya Terahara, Kenji Takizawa, Tayfun E. Tezduyar, Yuri Bazilevs, Ming‐Chen Hsu (2020). Heart valve isogeometric sequentially-coupled FSI analysis with the space–time topology change method. \href{https://doi.org/10.1007/s00466-019-01813-0}{DOI: 10.1007/s00466-019-01813-0}
\end{thebibliography}

\end{document}
